\subsection{Pregunta 6: auditoría de tablas de enrutamiento y redistribución}
\autor{Bárbara (Sage y distancias administrativas); Sebastián aporta la salida de Killjoy}

\subsubsection*{Killjoy}
\evid{Q6_tablas_rutas/show_ip_route_killjoy.png}{Salida de \cmd{show ip route} en Killjoy.}
\pend{señalar rutas O, O E2/E1 y B}

\subsubsection*{Sage}
\evid{Q6_tablas_rutas/show_ip_route_sage.png}{Salida de \cmd{show ip route} en Sage.}
\pend{señalar rutas B y las demás presentes}

\subsubsection*{Distancia administrativa}

La distancia administrativa (AD) es la confiabilidad que el router asigna a la fuente de una ruta; cuando dos protocolos aprenden el mismo prefijo, instala el de menor AD. Valores por defecto de Cisco IOS:

\begin{center}
\begin{tabular}{@{}lc@{}}
  \toprule
  Origen de la ruta & AD \\
  \midrule
  Red conectada (C) & 0 \\
  Ruta estática (S) & 1 \\
  eBGP (B) & 20 \\
  OSPF (O, O E1, O E2) & 110 \\
  iBGP (B) & 200 \\
  \bottomrule
\end{tabular}
\end{center}

\pend{comentar cuáles de estos valores aparecen entre corchetes en las capturas y qué efecto tienen en Killjoy, donde hay rutas OSPF y BGP}
